\documentclass{article}
\usepackage{thumbpdf,lmodern}
\usepackage{amsmath}
\usepackage{graphicx}
\usepackage{booktabs}
\usepackage[margin=1in]{geometry}

\newcommand{\pkg}[1]{\texttt{#1}}
\newcommand{\proglang}[1]{\textsf{#1}}
\newcommand{\code}[1]{\texttt{#1}}
\newcommand{\Plainauthor}[1]{}
\newcommand{\Plaintitle}[1]{}
\newcommand{\Shorttitle}[1]{}
\newcommand{\Plainkeywords}[1]{}
\newcommand{\Address}[1]{}
\newcommand{\Email}[1]{\texttt{#1}}
\newcommand{\myabstract}{}
\newcommand{\Abstract}[1]{\renewcommand{\myabstract}{#1}}
\newcommand{\mykeywords}{}
\newcommand{\Keywords}[1]{\renewcommand{\mykeywords}{#1}}
\newcommand{\Title}[1]{\title{#1}}

\author{Dylan A. Mordaunt\\Victoria University of Wellington}
\Title{Supplement B: Guide to Equity Weights}
\Abstract{
  This supplement provides a guide to the normative assumptions embedded in the Distributional Cost-Effectiveness Analysis (DCEA), specifically the derivation and interpretation of equity weights.
}
\Keywords{equity, DCEA, Atkinson index, ethics}

\begin{document}
\maketitle
\begin{abstract}
\myabstract
\end{abstract}

\section{The Normative Basis of Equity Weights}

The \pkg{vop\_poc\_nz} package uses the Atkinson social welfare function to operationalize equity.
This approach assumes that society has a preference for reducing health inequality, quantified by the inequality aversion parameter, $\epsilon$.

\subsection{Inequality Aversion Parameter ($\epsilon$)}

The parameter $\epsilon$ determines the curvature of the social welfare function.
\begin{itemize}
    \item $\epsilon = 0$: \textbf{Utilitarian}. A QALY is a QALY, regardless of who receives it. Equity weights are all 1.0.
    \item $\epsilon = 0.5$: \textbf{Mild Aversion}. Society is willing to sacrifice a small amount of total health to improve the health of the worst-off.
    \item $\epsilon = 1.0$: \textbf{Moderate Aversion}. The value of a health gain is inversely proportional to the recipient's current health status.
    \item $\epsilon > 10$: \textbf{Rawlsian (Maximin)}. Society prioritizes the worst-off group almost exclusively.
\end{itemize}

\section{Derivation of Weights}

The equity weight $w_g$ for subgroup $g$ is calculated as:
\begin{equation}
w_g = \left( \frac{H_{ref}}{H_g} \right)^\epsilon
\end{equation}
where $H_{ref}$ is the reference health level (typically the population mean or the best-off group).

Table~\ref{tab:weights} illustrates how weights change with $\epsilon$ for hypothetical subgroups.

\begin{table}[ht]
\centering
\begin{tabular}{lcccc}
\toprule
Subgroup & Baseline Health ($H_g$) & $\epsilon=0$ & $\epsilon=0.5$ & $\epsilon=1.0$ \\
\midrule
High SES (Ref) & 0.90 & 1.00 & 1.00 & 1.00 \\
Middle SES & 0.80 & 1.00 & 1.06 & 1.13 \\
Low SES & 0.70 & 1.00 & 1.13 & 1.29 \\
\bottomrule
\end{tabular}
\caption{Sensitivity of Equity Weights to Inequality Aversion ($H_{ref}=0.90$).}
\label{tab:weights}
\end{table}

\end{document}
